\subsection{逆行列normalizing包絡の全curvature判定}
\label{subsec:h3-order2-22-clifford-matrix}
\begin{theorem}\label{thm:h3-order2-22-clifford-matrix}
$F\in\mathrm{GL}(U)$、$D=\det F$、$h\in\operatorname{Sym}^2U$ とし、
\[
 B(F,h)=k\langle U,u,v\rangle/
 ([x,y],\ uT-F(T)u,\ vT-F^{-1}(T)v,\ vu+Duv-Dh),\quad T\in U
\]
とする。全係数でnoetherian quadratic Koszul AS4 domainであり、実際の中山は $++--$。
この全行列・全curvature族でH3切断を持つ包絡は $F=-I$ に限る。
その全切断・必要十分条件・全幾何学的矢は前節の $B_h$ の分類に一致する。
全rank4 graphの必要被覆はここでは主張しない。
\end{theorem}
\begin{proof}
$T_0=k[x,y]$ に $u$ を $F$ で加え、さらに
\[
 \sigma|_{T_0}=F^{-1},\quad \sigma(u)=-Du,\quad
 \delta|_{T_0}=0,\quad\delta(u)=Dh
\]
のOre変数 $v$ を加える。
$\delta$ が $uT=F(T)u$ を保つのは $hT=Th$ そのものである。
反復Oreの性質と四つのPBW overlapから主張した包絡性が従う。
Ore中山公式で正側は $F F^{-1}=I$、$u$ 側は $(-D)^{-1}D=-1$、
$v$ 側の先頭係数は $\det(F^{-1})(-D)=-1$。
$u\mapsto tu,v\mapsto t^{-1}v$ が一次補正を排除するので、中山は表示通りである。

正固有三次は十語 $f_3+A_1u^2+Luv+C_1v^2$ である。
正側との可換性は $A_1(F^2(T)-T)=0$ と
$C_1(F^{-2}(T)-T)=0$ を全 $T\in U$ に要求する。
$F^2\ne I$ なら $A_1=C_1=0$。
残る正常条件から、$u^2v$ と純正三次に $u$ を掛けた係数を比較して
\[
 F(L)=DL,\qquad (1+\operatorname{Sym}^3F)f_3=-D\,L F(h)
\]
を得る。$L=0$ なら $G$ は二元三次なので分解する。
$L\ne0$ なら $F$ は $U/kL$ に恒等として作用する。
上の第二式を $k[U]/(L)$ へ移すと $2f_3=0$、従って $L\mid f_3$。
よって $G=L(\text{二次式}+uv)$ は分解し、商はdomainにならない。
これは固有値や根の一般点だけの計算ではなく、全非involutionを同時に扱う。

$F^2=I$ は $I,-I$ とreflectionの三型だけである。
$F=I$ では全正常核は
\[
 G=L(uv-h/2),\qquad L\in U
\]
であり全て分解する。核の二列を直接検証でき、残る八列には定数minor256がある。
reflectionは $F(x)=x,F(y)=-y$ とできる。
$Gu+uG$ の $xu^3,xu^2v,xuv^2,x^3u$ 行は、
正常候補から $xu^2,xuv,xv^2,x^3$ の四係数を消去する。
この四列のminorは全 $h$ で16、残る六列はこの四行で零である。
従って全 $G$ が $y$ を因子に持ち、やはりdomainにならない。
$F=-I$ だけが残り、これは正確に前節の $B_h$ である。

全行列についての四overlapと、二つの例外involutionの全係数minorは
\[
\texttt{verify\_cubic\_order2\_22\_clifford\_matrix.py}
\]
で元の簡約規則から再生する。全H3の被覆や独立査読はこの検算の射程ではない。
\end{proof}
